\section{System and methodology}
\label{sec:method}

\subsection{Template-guided perception}
Each form template carries a QR payload identifying its template and version and four ArUco corner markers. The system decodes the QR, estimates a homography from detected corners, warps the image to a canonical canvas, and crops fields from versioned template coordinates. This converts a whole-page task into typed cell decisions and preserves the transformation parameters with the recognition attempt.

Digit crops are resized to $48\times64$ pixels. The built-in comparator binarizes them with Otsu's method and measures normalized mean absolute difference from ten rendered templates. If $d_1(x)$ and $d_2(x)$ are its best and second-best distances, its selection score is
\begin{equation}
  s(x)=d_2(x)-d_1(x).
\end{equation}
The evaluated whole-form configuration instead uses a HOG descriptor with a $48\times64$ window, $16\times16$ blocks, $8\times8$ cells and stride, and nine orientation bins, followed by a linear support-vector machine. Training uses 240 separately seeded rendered cells (eight seeds, three fonts, and ten digits). Its selection score is the difference between the largest and second-largest decision values. Calibration selected $\tau=0$, so all evaluation predictions were retained; the candidate--fact boundary still requires review before a prediction becomes authoritative. OMR cells use an interior fill ratio: at most 0.15 is unchecked, at least 0.35 is checked, and the interval between them is rejected for review.

\begin{figure}[t]
  \centering
  \resizebox{0.96\textwidth}{!}{\begin{tikzpicture}[node distance=5mm and 6mm,font=\scriptsize]
  \node[trustbox,minimum width=16mm] (capture) {\textbf{Capture}};
  \node[trustbox,minimum width=22mm,right=of capture] (locate) {\textbf{QR + ArUco}\\localization};
  \node[machine,minimum width=20mm,right=of locate] (recognize) {\textbf{Cell}\\recognition};
  \node[trustbox,minimum width=24mm,right=of recognize] (cert) {\textbf{Candidate}\\certificate};
  \node[trustbox,minimum width=24mm,right=of cert] (gate) {\textbf{Margin gate}\\$s(x)\geq\tau$?};
  \node[trustbox,minimum width=20mm,right=of gate] (queue) {\textbf{Review}\\queue};
  \node[human,minimum width=20mm,right=of queue] (review) {\textbf{Confirm or}\\correct};
  \node[fact,minimum width=18mm,right=of review] (export) {\textbf{Versioned}\\fact};
  \draw[flow] (capture)--(locate)--(recognize)--(cert)--(gate);
  \draw[flow] (gate)--node[above,font=\tiny]{candidate}(queue)--(review)--(export);
  \draw[flow] (gate.south) -- ++(0,-7mm) -| node[pos=.3,below,font=\tiny]{abstain} (queue.south);
\end{tikzpicture}
}
  \caption{Operational pipeline. Recognition outcomes cross the threshold gate into the review queue but remain non-authoritative until a confirm/correct action.}
  \label{fig:pipeline}
\end{figure}

\subsection{Candidate, decision, and fact services}
Evidence is assigned a SHA-256 digest. Recognition attempts and AI reviews are stored independently from confirmed record versions. The optional AI adapter accepts structured context and must return schema-valid review JSON; its prompt contract states that it cannot modify business facts. The default adapter is disabled, returning an explicit \texttt{NOT\_RUN} status. Local retrieval uses a dependency-free similarity index and exposes results only as references.

The review service accepts an expected record version, reviewer identifier, reason, candidate values, and evidence identifiers. It rejects a stale expected version, appends a new confirmed record on success, and emits a corresponding audit event. Exports query confirmed records only. Figure~\ref{fig:lineage} shows the reverse-trace chain.

\begin{figure}[t]
  \centering
  \resizebox{0.96\textwidth}{!}{\begin{tikzpicture}[node distance=8mm and 10mm,font=\small]
  \node[trustbox] (hash) {SHA-256 evidence};
  \node[trustbox,right=of hash] (attempt) {recognition attempt};
  \node[trustbox,right=of attempt] (candidate) {candidate ID};
  \node[human,right=of candidate] (event) {review event};
  \node[fact,right=of event] (version) {fact version};
  \node[fact,below=of version] (batch) {export batch};
  \draw[flow] (hash)--(attempt)--(candidate)--(event)--(version)--(batch);
  \draw[flow] (batch.west) -| node[pos=.25,below,font=\scriptsize]{reverse trace} (hash.south);
\end{tikzpicture}
}
  \caption{Evidence lineage from content-addressed input to export. Every arrow represents a persisted identifier used by the reverse-trace check.}
  \label{fig:lineage}
\end{figure}

\subsection{Implementation}
The evaluated artifact is a local single-machine prototype implemented with Python~3.11, Streamlit, SQLite, SQLAlchemy, Alembic, OpenCV, scikit-learn, and an optional provider adapter. Domain, application, adapter, and user-interface layers separate fact-transition services from machine-facing services. The experiment harness invokes these same repositories and application services for import, candidate persistence, confirmation, export, and trace queries.

The system retains local evidence, database, and export files as one backup unit. This local-first operational choice supports recovery in a small-plant setting \citep{kleppmann2019}, but it is not itself a security boundary. Deployment still requires host access control, encrypted backups, retention rules, and tested restore procedures.
